\chapter{Périmètre du projet}

\section{Périmètre inclus}

Le projet se concentre sur un MVP (Minimum Viable Product) permettant d'assurer les fonctionnalités essentielles pour l'achat et la vente de produits de seconde main entre étudiants.

Les fonctionnalités prioritaires incluses dans le MVP sont :

\begin{table}[h]
\centering
\renewcommand{\arraystretch}{1.25}
\begin{tabularx}{\textwidth}{>{\bfseries}c X c}
\toprule
\textbf{ID} & \textbf{Fonctionnalité} & \textbf{Priorité} \\
\midrule
F1 & Inscription et connexion des utilisateurs & Haute \\
F2 & Publication et gestion des annonces & Haute \\
F3 & Consultation du catalogue des produits & Haute \\
F4 & Recherche et filtrage des annonces & Haute \\
F5 & Consultation des détails d'un produit & Haute \\
F6 & Prise de contact entre acheteur et vendeur & Haute \\
F7 & Gestion de la disponibilité des produits (réservé / vendu) & Haute \\
F8 & Signalement et modération des annonces & Moyenne \\
F9 & Gestion du profil utilisateur & Moyenne \\
F10 & Consultation des annonces publiées par un utilisateur & Moyenne \\
\bottomrule
\end{tabularx}
\caption{Fonctionnalités incluses dans le MVP}
\end{table}

\section{Périmètre exclu}

Afin de rester réaliste compte tenu du temps disponible (six séances de TP), certaines fonctionnalités ne seront pas implémentées dans cette première version.

\begin{table}[h]
\centering
\renewcommand{\arraystretch}{1.25}
\begin{tabularx}{\textwidth}{>{\bfseries}X X}
\toprule
\textbf{Fonctionnalité exclue} & \textbf{Justification} \\
\midrule
Paiement en ligne & Nécessite une intégration de paiement et des mesures de sécurité supplémentaires. \\
Livraison intégrée et suivi des colis & Implique une logistique et des services externes qui dépassent le périmètre du prototype. \\
Application mobile native Android/iOS & Une application web responsive suffit pour la première version. \\
Système avancé de notation et de réputation & Peut être ajouté dans une version ultérieure après validation du fonctionnement principal. \\
Recommandations personnalisées par intelligence artificielle & Cette fonctionnalité n'est pas essentielle au besoin principal et augmenterait la complexité. \\
\bottomrule
\end{tabularx}
\caption{Fonctionnalités exclues du périmètre}
\end{table}

\section{Justification du périmètre}

La définition du périmètre permet de limiter la complexité, de mieux estimer les tâches et de concentrer les efforts de l'équipe sur les fonctionnalités prioritaires dans le temps imparti. \nomProjet\ vise ainsi à répondre au besoin réel des étudiants tout en restant réalisable dans le cadre de ce projet.
